\section{Agent-World: hacer que el entorno de entrenamiento y la política coevolucionen}

Dong Guanting (董冠霆) ubica, en Agent-World, el cuello de botella del self-evolving Agent en el entorno: Agentic RL necesita grandes cantidades de tareas ejecutables, verificables y con dificultad controlable, pero construir entornos de manera manual y etiquetar recompensas no puede mantener el ritmo del crecimiento de la política. La solución no consiste solo en ampliar el conjunto de datos, sino en construir una arena capaz de descubrir entornos de manera continua, sintetizar tareas, verificar resultados y diagnosticar brechas de capacidad.

\sourcefigure{0.94}{lecture04/slides-images/slide-002-00030s.jpg}{El objetivo de Agent-World: la inteligencia de la política y la oferta de entornos coevolucionan}{00:00:30--00:01:00}

\subsection{Por qué la insuficiencia de entornos es la principal restricción de Agentic RL}

El preentrenamiento de los modelos de lenguaje puede obtener una cantidad masiva de tokens a partir de corpus estáticos, pero el entrenamiento de Agent exige que cada muestra tenga estado, herramientas, acciones y retroalimentación ejecutable. Las API de servicios reales pueden ser inestables, costosas o involucrar privacidad; los sandbox construidos a mano, a su vez, tienen una cobertura limitada, lo que termina llevando a que la política se optimice una y otra vez sobre unos pocos benchmarks.

\sourcefigure{0.94}{lecture04/slides-images/slide-003-00060s.jpg}{Agentic Environment--Task Discovery: buscar entornos ejecutables en el espacio de temas}{00:01:00--00:04:15}

\begin{importantbox}{El entorno es parte de la distribución de entrenamiento}
El límite superior de la capacidad del Agent no lo determinan solo el modelo y el algoritmo, sino también el ecosistema de herramientas al que tiene acceso, los esquemas de bases de datos, la dificultad de las tareas y el verificador. Cuando el entorno es único, seguir aumentando los rollouts solo hará que se sobreajuste de manera más completa a un patrón de interacción limitado.
\end{importantbox}

\subsection{Primera etapa: Environment--Task Discovery}

Agent-World parte primero de temas reales a gran escala, en busca de bases de datos, API o conjuntos de herramientas ejecutables, y después los estandariza en una interfaz de interacción unificada. Enseguida, el sistema analiza el schema y las capacidades de las herramientas, y propone tareas que requieren consultas de varios pasos, combinaciones o razonamiento.

\sourcefigure{0.94}{lecture04/slides-images/slide-004-00255s.jpg}{El flujo de descubrimiento, de los temas del entorno y las capacidades de las herramientas a las tareas candidatas}{00:04:15--00:07:45}

\begin{knowledgebox}{Task discovery no es lo mismo que generar preguntas al azar}
Las tareas de alto valor deben depender del estado real dentro del entorno; la respuesta no se puede adivinar solo con el conocimiento previo del lenguaje. Las tareas también deben exigir varios pasos con herramientas y cubrir distintos modos de error. De lo contrario, el generador producirá pseudo tareas de Agent que parecen ricas pero que en realidad no requieren interacción.
\end{knowledgebox}

\subsection{Segunda etapa: Verifiable Task Synthesis}

El mayor riesgo de sintetizar tareas es no poder calificarlas de manera confiable. Agent-World construye el verifier mediante restricciones del entorno, ejecución de programas, consultas a bases de datos o revisión de reglas, de modo que el reward provenga en la medida de lo posible del estado externo y no del juicio subjetivo de otro modelo de lenguaje.

\sourcefigure{0.94}{lecture04/slides-images/slide-005-00465s.jpg}{Síntesis verificable de tareas: generar a la vez la tarea, la condición de ejecución y la lógica de verificación}{00:07:45--00:09:15}

Si el estado objetivo de la tarea $q$ es $g_q$ y la trayectoria del Agent es $\tau$, el verificador se puede escribir como

$$
r(\tau,q)=\mathbb{I}\!\left[V\!\left(s_{\mathrm{final}}(\tau),g_q\right)=1\right],
$$

donde $V$ es un programa de verificación determinista o aleatorio controlado. Para trayectorias largas, también se pueden agregar revisiones de formato, restricciones de herramientas y verificaciones del estado intermedio, para evitar que el Agent eluda la tarea real aprovechando una laguna de recompensa.

\begin{warningbox}{Que sea verificable no significa que no haya reward hacking}
Si el verifier solo revisa la cadena de texto final, el Agent puede falsificar la respuesta; si solo revisa el estado de la base de datos, puede lograr el objetivo mediante efectos secundarios no autorizados. El verificador debe cubrir a la vez el resultado, las restricciones del proceso y los límites de seguridad, y someter las tareas generadas a una auditoría independiente.
\end{warningbox}

\subsection{Tercera etapa: Continuous Self-Evolving Agent Training}

Una vez preparados el entorno y las tareas, el sistema entrena la política con Agentic RL en múltiples entornos, y después devuelve a la arena las debilidades expuestas en la evaluación. La arena no es un leaderboard fijo, sino un diagnosticador: identifica en qué combinaciones de herramientas, dependencias de largo alcance o recuperación de errores falla el modelo, y genera de manera dirigida la siguiente ronda de tareas.

\sourcefigure{0.94}{lecture04/slides-images/slide-008-00705s.jpg}{El ciclo cerrado de Agent-World: descubrimiento de entornos, síntesis de tareas, entrenamiento y diagnóstico}{00:11:45--00:12:30}

\sourcefigure{0.94}{lecture04/slides-images/slide-010-00840s.jpg}{El flujo de trabajo de Agentic RL en múltiples entornos y la arena de autoevolución}{00:14:00--00:15:00}

\begin{importantbox}{El significado de la coevolución}
Cuando la política se vuelve más fuerte, las tareas antiguas pierden su valor de entrenamiento; el generador de entornos debe aumentar la complejidad de las combinaciones o introducir nuevas herramientas. A la inversa, los fallos que exponen los nuevos entornos determinan el enfoque de la siguiente ronda de entrenamiento de la política. Solo cuando ambos lados cambian en conjunto se puede hablar propiamente de agent--world co-evolution.
\end{importantbox}

\subsection{Resultados experimentales e interpretación correcta}

El informe compara, en varios benchmarks de Agent, el modelo base, la sola adición de datos de entorno y el ciclo cerrado completo de Agent-World. La tendencia general es que el entrenamiento a través de múltiples entornos ofrece una mejor transferencia promedio que el entrenamiento sobre un solo benchmark, y las tareas generadas de manera dirigida también mejoran debilidades específicas.

\sourcefigure{0.94}{lecture04/slides-images/slide-012-00960s.jpg}{Los principales resultados de Agent-World en varios benchmarks}{00:16:00--00:18:00}

\sourcefigure{0.94}{lecture04/slides-images/slide-015-01215s.jpg}{El cambio en la cobertura de capacidades que trae la expansión de entornos y tareas}{00:20:15--00:21:00}

\sourcefigure{0.94}{lecture04/slides-images/slide-018-01320s.jpg}{Ablación de los distintos componentes y ganancias de entrenamiento}{00:22:00--00:22:30}

\begin{warningbox}{Que el benchmark mejore no significa que el mundo real ya esté resuelto}
El entorno sintético todavía puede ser más ordenado que un sistema real: los errores de herramientas, los cambios de permisos, la latencia y la falta de determinismo son todos más débiles. El experimento respalda que «la síntesis escalable de entornos es efectiva», pero de ahí no se puede concluir directamente que el Agent ya sea capaz de actuar de manera autónoma y prolongada en la internet abierta o en sistemas empresariales.
\end{warningbox}

\subsection{Resumen de la sección}

Agent-World eleva el problema de los datos de entrenamiento a un problema de ingeniería del entorno. El sistema primero descubre entornos ejecutables, después sintetiza tareas con verifier, entrena la política mediante RL en múltiples entornos, y usa la arena para diagnosticar debilidades y generar el currículo de la siguiente ronda. Esto se conecta de manera natural con los dos capítulos anteriores: el causal world model ayuda a entender el entorno, la rollout infrastructure se encarga de la interacción a escala, y Agent-World provee de manera continua nuevos mundos con los que vale la pena interactuar.
